% !TEX program = lualatex
\documentclass[a4paper,11pt]{ltjsarticle}
\usepackage{luatexja-fontspec}
\usepackage{amsmath,amssymb}
\usepackage[margin=18mm,landscape]{geometry}
\usepackage{graphicx}
\setlength{\parindent}{0pt}
\title{$^{151}$Eu ($F=6$) 量子化軸回転行列 $R_a(\theta)=e^{-i\theta F_a}$ の成分表示}
\author{KozumaLab / mitsuki}
\date{\today}
\begin{document}
\maketitle

\section*{1. 定義と $x\leftrightarrow y$ の関係}
軸 $a$ まわりの回転 $R_a(\theta)=e^{-i\theta F_a}$（$\hbar=1$）。$F=6$ で $13\times13$。
$F_y$ は純虚・反対称なので $R_y(\theta)$ は\textbf{実行列}＝Wigner の $d$ 行列：
\[
[R_y(\theta)]_{m'm}=d^{6}_{m'm}(\theta)\in\mathbb{R},\qquad
[R_x(\theta)]_{m'm}=i^{\,m'-m}\,d^{6}_{m'm}(\theta).
\]
（後者は数値で誤差 $10^{-16}$ で確認。すなわち $|R_x|=|R_y|=d$、違いは位相 $i^{m'-m}$ のみ。$R_x$ は $R_y$ を $R_z(\pm\pi/2)$ で挟んだもの。）
$R_y$ は直交（$R R^{\mathsf T}=I$、残差 $3.3\times10^{-16}$）、$\det=+1$。

\section*{2. 一般角 $\theta$ の閉形式（Wigner の $d$ 公式）}
\[
d^{j}_{m'm}(\theta)=\!\sum_{k}\!\frac{(-1)^{k}\sqrt{(j{+}m')!(j{-}m')!(j{+}m)!(j{-}m)!}}
{(j{+}m{-}k)!\,k!\,(m'{-}m{+}k)!\,(j{-}m'{-}k)!}
\Bigl(\cos\tfrac\theta2\Bigr)^{2j+m-m'-2k}\!\Bigl(\sin\tfrac\theta2\Bigr)^{m'-m+2k}
\]
（$j=6$、$k$ は全階乗が非負の範囲で和。）

\section*{3. 数値：$d^{6}(\theta{=}16^\circ)=R_y(16^\circ)$（実、成分表示）}
\begin{center}\resizebox{\textwidth}{!}{$\begin{array}{c|ccccccccccccc} & m'{=}6 & m'{=}5 & m'{=}4 & m'{=}3 & m'{=}2 & m'{=}1 & m'{=}0 & m'{=}-1 & m'{=}-2 & m'{=}-3 & m'{=}-4 & m'{=}-5 & m'{=}-6\\ \hline m{=}6 & 0.8893 & -0.4329 & 0.1427 & -0.0366 & 0.0077 & -0.0014 & 0.0002 & 0 & 0 & 0 & 0 & 0 & 0\\ m{=}5 & 0.4329 & 0.6961 & -0.5283 & 0.2123 & -0.0609 & 0.0137 & -0.0025 & 0.0004 & 0 & 0 & 0 & 0 & 0\\ m{=}4 & 0.1427 & 0.5283 & 0.5536 & -0.5661 & 0.2587 & -0.0797 & 0.0186 & -0.0034 & 0.0005 & -0.0001 & 0 & 0 & 0\\ m{=}3 & 0.0366 & 0.2123 & 0.5661 & 0.4519 & -0.5802 & 0.2885 & -0.0924 & 0.0217 & -0.0039 & 0.0006 & -0.0001 & 0 & 0\\ m{=}2 & 0.0077 & 0.0609 & 0.2587 & 0.5802 & 0.3840 & -0.5846 & 0.3052 & -0.0987 & 0.0228 & -0.0039 & 0.0005 & 0 & 0\\ m{=}1 & 0.0014 & 0.0137 & 0.0797 & 0.2885 & 0.5846 & 0.3450 & -0.5855 & 0.3105 & -0.0987 & 0.0217 & -0.0034 & 0.0004 & 0\\ m{=}0 & 0.0002 & 0.0025 & 0.0186 & 0.0924 & 0.3052 & 0.5855 & 0.3323 & -0.5855 & 0.3052 & -0.0924 & 0.0186 & -0.0025 & 0.0002\\ m{=}-1 & 0 & 0.0004 & 0.0034 & 0.0217 & 0.0987 & 0.3105 & 0.5855 & 0.3450 & -0.5846 & 0.2885 & -0.0797 & 0.0137 & -0.0014\\ m{=}-2 & 0 & 0 & 0.0005 & 0.0039 & 0.0228 & 0.0987 & 0.3052 & 0.5846 & 0.3840 & -0.5802 & 0.2587 & -0.0609 & 0.0077\\ m{=}-3 & 0 & 0 & 0.0001 & 0.0006 & 0.0039 & 0.0217 & 0.0924 & 0.2885 & 0.5802 & 0.4519 & -0.5661 & 0.2123 & -0.0366\\ m{=}-4 & 0 & 0 & 0 & 0.0001 & 0.0005 & 0.0034 & 0.0186 & 0.0797 & 0.2587 & 0.5661 & 0.5536 & -0.5283 & 0.1427\\ m{=}-5 & 0 & 0 & 0 & 0 & 0 & 0.0004 & 0.0025 & 0.0137 & 0.0609 & 0.2123 & 0.5283 & 0.6961 & -0.4329\\ m{=}-6 & 0 & 0 & 0 & 0 & 0 & 0 & 0.0002 & 0.0014 & 0.0077 & 0.0366 & 0.1427 & 0.4329 & 0.8893\end{array}$}\end{center}
$R_x(16^\circ)$ は各成分に位相 $i^{m'-m}$ を掛けたもの（大きさは上表と同一）。

\section*{4. 物理の中心：旧 $m{=}{-}6$ が $\theta$ 傾けで漏れる先（$-6$ ソース列）}
$c\equiv\cos\tfrac\theta2,\ s\equiv\sin\tfrac\theta2$ として単項に縮約：
\[
d^{6}_{m,-6}(\theta)=(-1)^{6-m}\sqrt{\tbinom{12}{6+m}}\;c^{\,6-m}s^{\,6+m}.
\]
\[
R_y(\theta)\,|{-}6\rangle=\bigl(
s^{12},\;
-\sqrt{12}\,cs^{11},\;
\sqrt{66}\,c^{2}s^{10},\;
-\sqrt{220}\,c^{3}s^{9},\;
\sqrt{495}\,c^{4}s^{8},\;
-\sqrt{792}\,c^{5}s^{7},\;
\sqrt{924}\,c^{6}s^{6},
\]
\[
-\sqrt{792}\,c^{7}s^{5},\;
\sqrt{495}\,c^{8}s^{4},\;
-\sqrt{220}\,c^{9}s^{3},\;
\sqrt{66}\,c^{10}s^{2},\;
-\sqrt{12}\,c^{11}s,\;
c^{12}
\bigr)^{\mathsf T}\quad(m=6,\dots,-6).
\]
占有数（移行確率）は\textbf{二項分布}になる：
\[
\bigl|d^{6}_{m,-6}(\theta)\bigr|^{2}
=\binom{12}{6+m}\bigl(\sin^{2}\tfrac\theta2\bigr)^{6+m}\bigl(\cos^{2}\tfrac\theta2\bigr)^{6-m}
=\mathrm{Binom}\!\Bigl(12,\,6{+}m;\,p{=}\sin^{2}\tfrac\theta2\Bigr).
\]
例：$\theta{=}16^\circ\Rightarrow p{=}\sin^2 8^\circ{=}0.0194$、$P({-}6){=}(1{-}p)^{12}{=}0.790$、$P({-}5){=}12p(1{-}p)^{11}{=}0.187$。
$\theta{=}90^\circ\Rightarrow p{=}1/2$（対称二項、最大混合）。
\end{document}
